17.0 WARUNKI ZWYCIĘSTWA

Ogólna zasada:

Istnieją dwa rodzaje zwycięstwa: zwycięstwo zespołowe i zwycięstwo indywidualne.
Grupa, tzn. oddział wychodzi z gry zwycięsko jeżeli odnjdzie Czarne Wrota, zniszczy je i co najmniej jedna postać wyjdzie żywa z labiryntu.
Nie wymagane jest rozgrywanie więcej niż jednej przygody tzn. więcej niż jednej wizyty w labiryncie.

Zwycięzcą indywidualnym zostaje ta postać, która ujdzie z życiem ze wszystkich rozegranych przygód, zdobędzie najwięcej punktów doświadczenia oraz zgromadzi największy majątek w sztukach złota i kosztownościach.
Jeżeli gracze decydują się na zakończenie gry przed znalezieniem lub zniszczeniem Czarnych Wrót zwycięstwo grupowe jest niemożliwe.

Sposób postępowania:

Oddział wchodzi do labiryntu i poszukuje Czarnych Wrót.
Jeżeli w jakimś momencie gracze decydują, że istnieje zbyt duże prawdopodobieństwo, że oddział zostanie wybity do nogi zanim je odnajdzie, może on wyjść z labiryntu tą samą drogą, którą do niego wszedł, tzn. przez bramę labiryntu na pierwszym poziomie.
Taki fragment gry – od chwili wejścia do labiryntu do chwili opuszczenia go bez wykonania głównego zadania – nazywany jest przygodą.
Po wyleczeniu wszystkich ran oddział może ponownie wejść do labiryntu.
Zbadane podczas poprzednich przygód segmenty są znane w przygodach następnych.

Zasady szczegółowe:

[17.1] Sztuki złota i kosztowności, zdobyte przez postać w trakcie jednej przygody są zostawiane w domu przed rozpoczęciem następnej, o ile nie zostały wydane na zwiększenie poszczególnych współczynników (patrz 12.2).

[17.2] Potwory, opanowane przy pomocy czaru ,,Urok'' zostają wyzwolone spod jego władzy w momencie wyjścia oddziału z labiryntu.

[17.3] Po każdej przygodzie, w trakcie której zostanie zabity jeden (lub więcej) członek oddziału można wprowadzić do gry nowe postacie, tak by oddział liczył 3 bohaterów i 3 uczniów.
Nowo wprowadzone postacie zaczynają oczywiście z zerowym kontem punktów doświadczenia i skarbów.

[17.4] Przed każdym kolejnym wejściem do labiryntu postać może wybrać nową broń, jednak tak, by miała przy sobie nie więcej niż dwa jej rodzaje.

[17.5] Zakłada się, że między wyjściem z labiryntu a ponownym do niego wejściem upływa 6 miesięcy.
W ciągu tego czasu postacie nie tylko leczą swe zranienia (odcinające rany), ale takze ,,starzeja się'', a więc wolniej nabywają nowe, wyższe umiejętności.
Jest to odzwierciedlone w rosnącym koszcie nabywania punktów, zwiększających poszczególne współczynniki.
Koszty te rosną o 50 po zakończeniu każdej kolejnej przygody, poczynając od drugiej (tzn. od drugiego wyjścia z labiryntu).

4.
Gracze wybierają dwa rodzaje broni dla każdego ze swoich uczniów (dowolnie z tabeli 9.9).

5.
Gracze rzucają 1K6 i z tabeli 4.4 ,,Potencjał magiczny'' odczytują wielkość współczynnika potencjał magiczny dla każdego ucznia, a następnie zapisują go na karcie cech.
Każdy współczynnik składa się z trzech cyfr.

6.
Dla każdej postaci, której potencjał magiczny (jeden z jego trzech członów) jest większy od zera gracze wybierają czary, którymi ta postać będzie dysponować w czasie gry.
Liczba czarów musi być równa największemu z trzech członów, składających się na potencjał magiczny.
Na przykład: postać, posiadająca potencjał magiczny równy 4/3/2 powinna rozpocząć grę znając 4 różne czary.

7.
Zakłada się, że każdy uczeń zdobył już w swoim dotychczasowym życiu pewne doświadczenie.
W związku z tym gracz może:

a. zwiększyć jego siłę o 1 lub b. zwiększyć jego biegłość w jego podstawowej broni o 1 lub c. zwiększyć jego zdolności o 1.

C. USTALENIE SZYKU MARSZOWEGO.

[Patrz opis w rozdziale 4.0 - Przygotowanie do gry, punkt C]

Szyk marszowy składa się z pewnej liczby szeregów, przy czym w żadnym z nich nie może być więcej niż trzy postacie.
Gracze układają żetony, oznaczające postacie w dogodnym miejscu na stole, tak aby tworzyły one szyk marszowy.
Od decyzji graczy zależy która postać będzie w jakim szeregu i na jakiej pozycji.
W pierwszym szeregu szyku muszą być co najmniej dwie postacie (o ile aż tyle pozostaje przy życiu).

3.
Gracze mogą dowolnie zmieniać szyk marszowy do chwili rozpoczęcia walki.
W jej trakcie szyk może zostać zmieniony tylko podczas fazy ,,reorganizacja oddziału'' (patrz 9.8).

D.
Ustalenie dominującego słońca.

Jeden z graczy rzuca 1K6 w celu ustalenia które z trzech słońc jest dominujące czyli pod wpływem którego słońca oddział się teraz znajduje.

Wpływ słońca określa liczbę czarów, które mogą być używane przez daną postać podczas bieżącej przygody (czyli w ciągu części gry liczonej od momentu wejścia do labiryntu do momentu wyjścia z niego w celu wyleczenia ran, zmiany broni, itp.).
Dla każdej kolejnej przygody należy na nowo ustalać dominujące słońce.
Rezultaty rzutu kostką odczytuje się następująco: